\section{Программная реализация задачи}

\subsection{Структура программы}

Программа реализована на языке Java 17 с графическим интерфейсом на Swing. Структура проекта:

\begin{itemize}
\item \texttt{ode/} — интерфейс дифференциального уравнения и три реализации ($y' = y$, $y' = x + y$, $y' = -2xy$)
\item \texttt{methods/} — численные методы (модифицированный метод Эйлера, метод Рунге-Кутта 4-го порядка, метод Милна)
\item \texttt{model/} — модель данных (SolveResult)
\item \texttt{gui/} — графический интерфейс (InputPanel, ResultTablePanel, PlotPanel, MainFrame)
\end{itemize}

Графический интерфейс разделён на две части:
\begin{itemize}
\item Верхняя часть — выбор уравнения, ввод начальных параметров ($x_0$, $y_0$, $x_n$, $h$, $\varepsilon$), выбор методов, кнопка <<Вычислить>>
\item Нижняя часть — левая панель с таблицей результатов, правая панель с графиком
\end{itemize}

\subsection{Интерфейс дифференциального уравнения}

\begin{lstlisting}[style=javaStyle, label=code:ode]
public interface ODE {
    double f(double x, double y);
    double exact(double x, double x0, double y0);
    String expression();
}
\end{lstlisting}

\subsection{Модифицированный метод Эйлера}

\begin{lstlisting}[style=javaStyle, label=code:euler]
public class ImprovedEulerMethod {
    private static final int ORDER = 2;

    public static double step(ODE ode, double xi, double yi, double h) {
        double f_xi_yi = ode.f(xi, yi);
        double yTilde = yi + h * f_xi_yi;
        return yi + (h / 2.0) * (f_xi_yi + ode.f(xi + h, yTilde));
    }

    public static SolveResult solve(ODE ode, double x0, double y0,
                                     double xn, double h) {
        int n = (int) Math.round((xn - x0) / h);
        double[] x = new double[n + 1];
        double[] y = new double[n + 1];
        x[0] = x0;
        y[0] = y0;
        for (int i = 0; i < n; i++) {
            x[i + 1] = x[i] + h;
            y[i + 1] = step(ode, x[i], y[i], h);
        }
        double rungeErr = rungeError(ode, x0, y0, xn, h);
        return new SolveResult(x, y, rungeErr, "Modified Euler", "runge");
    }

    public static double rungeError(ODE ode, double x0, double y0,
                                     double xn, double h) {
        int n = (int) Math.round((xn - x0) / h);
        SolveResult half = solveHalfStep(ode, x0, y0, xn, h);
        double maxErr = 0.0;
        double yi_h = y0, xi = x0;
        for (int i = 0; i <= n; i++) {
            double diff = Math.abs(yi_h - half.getY()[2 * i]);
            double r = diff / (Math.pow(2, ORDER) - 1);
            if (r > maxErr) maxErr = r;
            if (i < n) {
                yi_h = step(ode, xi, yi_h, h);
                xi += h;
            }
        }
        return maxErr;
    }
}
\end{lstlisting}

\subsection{Метод Рунге-Кутта четвёртого порядка}

\begin{lstlisting}[style=javaStyle, label=code:rk4]
public class RungeKuttaMethod {
    private static final int ORDER = 4;

    public static double step(ODE ode, double xi, double yi, double h) {
        double k1 = h * ode.f(xi, yi);
        double k2 = h * ode.f(xi + h / 2.0, yi + k1 / 2.0);
        double k3 = h * ode.f(xi + h / 2.0, yi + k2 / 2.0);
        double k4 = h * ode.f(xi + h, yi + k3);
        return yi + (k1 + 2.0 * k2 + 2.0 * k3 + k4) / 6.0;
    }

    public static SolveResult solve(ODE ode, double x0, double y0,
                                     double xn, double h) {
        int n = (int) Math.round((xn - x0) / h);
        double[] x = new double[n + 1];
        double[] y = new double[n + 1];
        x[0] = x0;
        y[0] = y0;
        for (int i = 0; i < n; i++) {
            x[i + 1] = x[i] + h;
            y[i + 1] = step(ode, x[i], y[i], h);
        }
        double rungeErr = rungeError(ode, x0, y0, xn, h);
        return new SolveResult(x, y, rungeErr, "Runge-Kutta 4", "runge");
    }
}
\end{lstlisting}

\subsection{Метод Милна}

\begin{lstlisting}[style=javaStyle, label=code:milne]
public class MilneMethod {

    public static SolveResult solve(ODE ode, double x0, double y0,
                                     double xn, double h) {
        int n = (int) Math.round((xn - x0) / h);
        double[] x = new double[n + 1];
        double[] y = new double[n + 1];
        double[] f = new double[n + 1];

        x[0] = x0;
        y[0] = y0;
        f[0] = ode.f(x0, y0);

        // \u0417\u0430\u043f\u0443\u0441\u043a: \u0432\u044b\u0447\u0438\u0441\u043b\u0435\u043d\u0438\u0435 y1, y2, y3 \u043c\u0435\u0442\u043e\u0434\u043e\u043c \u0420\u0443\u043d\u0433\u0435-\u041a\u0443\u0442\u0442\u0430
        for (int i = 0; i < Math.min(3, n); i++) {
            x[i + 1] = x[i] + h;
            y[i + 1] = RungeKuttaMethod.step(ode, x[i], y[i], h);
            f[i + 1] = ode.f(x[i + 1], y[i + 1]);
        }

        // \u041f\u0440\u0435\u0434\u0438\u043a\u0442\u043e\u0440-\u043a\u043e\u0440\u0440\u0435\u043a\u0442\u043e\u0440 \u041c\u0438\u043b\u043d\u0430
        for (int i = 4; i <= n; i++) {
            x[i] = x[i - 1] + h;

            // \u041f\u0440\u0435\u0434\u0438\u043a\u0442\u043e\u0440
            double yPredict = y[i - 4]
                + (4.0 * h / 3.0) * (2.0 * f[i - 3] - f[i - 2] + 2.0 * f[i - 1]);

            // \u041a\u043e\u0440\u0440\u0435\u043a\u0442\u043e\u0440
            double fPredict = ode.f(x[i], yPredict);
            y[i] = y[i - 2] + (h / 3.0) * (f[i - 2] + 4.0 * f[i - 1] + fPredict);
            f[i] = ode.f(x[i], y[i]);
        }

        // \u041e\u0446\u0435\u043d\u043a\u0430 \u043f\u043e\u0433\u0440\u0435\u0448\u043d\u043e\u0441\u0442\u0438 \u043f\u043e \u0442\u043e\u0447\u043d\u043e\u043c\u0443 \u0440\u0435\u0448\u0435\u043d\u0438\u044e
        double maxErr = 0.0;
        for (int i = 0; i <= n; i++) {
            double exact = ode.exact(x[i], x0, y0);
            double diff = Math.abs(exact - y[i]);
            if (diff > maxErr) maxErr = diff;
        }
        return new SolveResult(x, y, maxErr, "Milne", "exact");
    }
}
\end{lstlisting}

\subsection{Графический интерфейс}

Приложение реализует следующие возможности:
\begin{enumerate}
\item Выбор одного из трёх предустановленных дифференциальных уравнений
\item Ручной ввод начальных данных: $x_0$, $y_0$, $x_n$, $h$, $\varepsilon$
\item Выбор реализованных методов (модифицированный метод Эйлера, метод Рунге-Кутта 4, метод Милна)
\item Отображение результатов в виде таблицы (значения $x_i$, $y_i$ для каждого метода, точное решение)
\item Отображение погрешностей (правило Рунге для одношаговых методов, сравнение с точным решением для метода Милна)
\item Построение графика точного решения и приближённых решений
\item Обработка некорректных входных данных ($x_n \leq x_0$, $h \leq 0$, нечисловые значения, не выбран ни один метод)
\end{enumerate}

\section{Результаты работы программы}

\subsection{Тестовый набор 1}

Уравнение $y' = y$, начальное условие $y(0) = 1$, интервал $[0, 2]$, шаг $h = 0{,}2$, точность $\varepsilon = 0{,}001$.

Здесь размещается скриншот работы программы для тестового набора 1.

\subsection{Тестовый набор 2}

Уравнение $y' = x + y$, начальное условие $y(0) = 1$, интервал $[0, 2]$, шаг $h = 0{,}2$, точность $\varepsilon = 0{,}001$.

Здесь размещается скриншот работы программы для тестового набора 2.

\subsection{Тестовый набор 3}

Уравнение $y' = -2xy$, начальное условие $y(0) = 1$, интервал $[0, 1]$, шаг $h = 0{,}1$, точность $\varepsilon = 0{,}001$.

Здесь размещается скриншот работы программы для тестового набора 3.

\section{Графики точного решения и приближённых решений}

\subsection{График для уравнения $y' = y$}

Здесь размещается график точного решения и приближённых решений для уравнения $y' = y$.

\subsection{График для уравнения $y' = x + y$}

Здесь размещается график точного решения и приближённых решений для уравнения $y' = x + y$.

\subsection{График для уравнения $y' = -2xy$}

Здесь размещается график точного решения и приближённых решений для уравнения $y' = -2xy$.

\section{Вывод}

В ходе выполнения лабораторной работы были изучены и реализованы три численных метода решения задачи Коши для обыкновенных дифференциальных уравнений первого порядка: модифицированный метод Эйлера (одношаговый, 2-й порядок точности), метод Рунге-Кутта четвёртого порядка (одношаговый, 4-й порядок точности) и метод Милна (многошаговый метод прогноза и коррекции, 4-й порядок точности).

Для одношаговых методов (модифицированный метод Эйлера и метод Рунге-Кутта) оценка погрешности выполнена по правилу Рунге. Для многошагового метода Милна погрешность оценена путём сравнения с точным решением: $\varepsilon = \max |y_{i \text{ точн}} - y_i|$.

Программная реализация выполнена на языке Java с графическим интерфейсом Swing. Приложение позволяет выбирать уравнение из трёх предложенных, задавать начальные параметры, выбирать реализованные методы, получать таблицу приближённых значений и график сравнения численных решений с точным.

Теоретические оценки точности методов подтверждены вычислительными экспериментами: метод Рунге-Кутта 4-го порядка показал наилучшую точность среди реализованных методов, метод Милна продемонстрировал сопоставимую точность при меньших вычислительных затратах на каждом шаге за счёт использования значений, вычисленных на предыдущих шагах.
